% Shared preamble for the printable cards.
%
% Each card is a standalone document that must come out as ONE A4 side. That
% constraint is the whole design: a card is consulted under pressure, and a
% card that runs to two sides is a document, which is not what a desk needs.
%
% tools/check_links.py does not police this. The build does: check-cards in the
% Makefile fails any card whose PDF has more than one page.
\documentclass[10pt,a4paper]{article}
\usepackage[T1]{fontenc}
\usepackage[utf8]{inputenc}
\IfFileExists{lmodern.sty}{\usepackage{lmodern}}{}
\usepackage[a4paper,margin=14mm]{geometry}
\usepackage{xcolor}
\usepackage{enumitem}
\usepackage{tabularx}
\usepackage{booktabs}
\usepackage{titlesec}

\definecolor{kitink}{RGB}{31,78,121}
\definecolor{kitwarn}{RGB}{150,30,30}
\definecolor{kitgrey}{RGB}{110,110,110}

\pagestyle{empty}
\setlength{\parindent}{0pt}
\setlength{\parskip}{3pt}
\setlist{itemsep=1pt, topsep=2pt, parsep=0pt, leftmargin=*}

\titleformat{\section}{\bfseries\color{kitink}}{}{0pt}{}
\titlespacing*{\section}{0pt}{7pt}{2pt}

\newcommand{\cardtitle}[2]{%
  {\LARGE\bfseries\color{kitink} #1\par}
  {\color{kitgrey}\small #2\par}
  \vspace{2pt}{\color{kitink}\rule{\linewidth}{1pt}}\par\vspace{4pt}}

\newcommand{\cardfoot}[1]{%
  \vfill{\color{kitink}\rule{\linewidth}{0.4pt}}\par\vspace{2pt}
  {\scriptsize\color{kitgrey} #1\par}}

\newcommand{\warn}[1]{\textcolor{kitwarn}{\textbf{#1}}}
% Same spelling as the book's preamble. Defined again rather than shared,
% because a card must build standalone -- somebody printing one should not
% need the book's preamble on the path.
\newcommand{\dash}{\,---\,}

\begin{document}
\cardtitle{Banned claims}{Fill in per process. Read immediately before every call.}

The sentences I must not say: convenient, half-believed, and not supported by
my own evidence.

\section{The list}
\small
\begin{tabularx}{\linewidth}{@{}XXX@{}}
\toprule
\textbf{Must not say} & \textbf{Why my evidence does not support it} & \textbf{What I say instead} \\
\midrule
\rule{0pt}{16pt} & & \\
\rule{0pt}{16pt} & & \\
\rule{0pt}{16pt} & & \\
\rule{0pt}{16pt} & & \\
\bottomrule
\end{tabularx}
\normalsize

\section{Check yourself against these shapes}
\begin{itemize}
  \item A capability my integration has \textbf{never exercised live} \dash\ I
        built it against a specification.
  \item A \textbf{count} I cannot reproduce today: tests, scenarios, how much
        something improved.
  \item An outcome I \textbf{contributed to}, said as though I owned it.
  \item A \warn{superlative} \dash\ only, first, every time, always. The most
        dangerous entries: they need one counterexample and interviewers often
        have one.
  \item Work by somebody else that I was near.
\end{itemize}

\section{The one that would cost most}
If I say one thing on this list, which one ends the conversation?
\par\vspace{4pt}\rule{\linewidth}{0.4pt}\par\vspace{7pt}

\section{The sentence that is stronger than the one I was going to say}
{\color{kitgrey}If a deliberately broken variant survived my own checks, say
that. It is the story of an instrument being tested and found imperfect, which
is what rigour looks like from outside \dash\ and it is stronger than
``it caught every one'', which is usually false by one case.}

\cardfoot{Staff AI Engineer Preparation Kit -- Chapter 10.}
\end{document}
